% !TEX root = print.tex
% In this file you should put the actual content of the blueprint.
% It will be used both by the web and the print version.
% It should *not* include the \begin{document}
%
% If you want to split the blueprint content into several files then
% the current file can be a simple sequence of \input. Otherwise It
% can start with a \section or \chapter for instance.

This project formalizes the theory of Knight Models, which describe a construction of Julia Knight in order to find a $\mathcal{L}_{\omega_1, \omega}$-sentence that characterizes $\aleph_1$ in the sense that there is a model of the sentence of cardinality $\aleph_1$ but no larger \cite{Knight1977}.
Greg Hjorth realized that their automorphism groups have interesting properties from the perspective of invariant descriptive set theory \cite{Hjorth1999} \cite{Hjorth2002Knight} \cite{Hjorth2010}.

We axiomatize the Knight models, the blueprints (which essentially the $\mathcal{L}_{\infty, \omega}$-types of the Knight models) and the KnightAutGroups (which are the automorphism groups of the Knight models).
In constructing a Knight model, we first construct a KnightAutGroup, then we use it to construct a blueprint, which we then use to construct a Knight model.
Finally, a Zorn's lemma argument is used to construct a Knight model of cardinality $\aleph_1$.

We first start with the definition of a Knight model, which in our formalization depends on the definition of a blueprint.

\begin{definition}[Blueprint]
  \label{def:blueprint}
  \lean{Blueprint}
  \leanok
  A \emph{blueprint} $(\Omega, {\le}, {\cdot})$ is a countable right-cancellative monoid and a dense linear order with no minimum satisfying additionally
  \begin{enumerate}
  \item $1$ is the maximum element;
  \item $\forall a, b, \; a * b \le b$;
  \item if $a \le b$ then for some $c$,\ $a = c \cdot b$;
  \item if $a, b < 1$ then for some $c$,\ $a = b \cdot c$.
  \end{enumerate}
\end{definition}

Note that the existence of a blueprint is not obvious and seems to require a Baire category argument.

\begin{definition}[Knight model]
  \label{def:knight-model}
  \lean{KnightModel}
  \leanok
  \uses{def:blueprint}
  A Knight model with blueprint $(\Omega, {\le}, {\circ})$ is a structure $\mathcal{K} = (K, {\preceq}, f_n)_{n \in \Omega}$ such that $\preceq$ is a dense linear order on $K$ without endpoints satisfying
  \begin{enumerate}
  \item $\forall a \preceq b \in K,\ \exists n \in \Omega,\ f_n(b) = a$;
  \item $\forall a \in K,\ \forall m, n \in \Omega,\ f_m(f_n(a)) = f_{m \circ n}(a)$;
  \item $\forall a \in K,\ \forall m, n \in \Omega,\ f_m(a) \preceq f_n(a) \leftrightarrow m \le n$.
  \end{enumerate}
\end{definition}

Here's the main theorem.

\begin{theorem}[Main theorem]
  \label{thm:main-theorem}
  \uses{thm:knight-model-card-bound, thm:chain-uncountable-model}
  A blueprint exists, and for any blueprint $\Omega$ there are Knight models with blueprint $\Omega$ of cardinality $\aleph_0$ and $\aleph_1$ but no other cardinalities, and moreover any two Knight models with blueprint $\Omega$ of cardinality $\aleph_0$ are isomorphic.
\end{theorem}

\begin{proof}
  \leanok
  In Section~\ref{sec:knight-aut-group} we define a KnightAutGroup and show one exists, Theorem \ref{thm:knight-aut-group}.
  In Section~\ref{sec:blueprint} we derive a blueprint from the KnightAutGroup, Theorem \ref{thm:knight-aut-group-is-blueprint}.

  In Section~\ref{sec:knight-model} we derive a countable Knight model from a blueprint, Theorem \ref{thm:blueprint-is-knight-model}.
  Uniqueness of countable Knight models of the same blueprint up to isomorphism is Theorem \ref{thm:countable-generic-iso}.
  The upper bound of $\aleph_1$ is Theorem~\ref{thm:knight-model-card-bound}.

  In Section~\ref{sec:chains} we set up the machinery to use Zorn's lemma to construct a Knight model of cardinality $\aleph_1$, Theorem~\ref{thm:chain-uncountable-model}.
\end{proof}



\section{KnightAutGroup}\label{sec:knight-aut-group}

\begin{definition}[KnightAutGroup]
  \label{def:knight-aut-group}
  \lean{KnightAutGroup}
  \leanok
  A \emph{KnightAutGroup} is a subgroup $H \le \mathrm{Aut}(\mathbb{Q}, {<})$ satisfying:
  \begin{enumerate}
  \item $\forall a, b \in \mathbb{Q},\ \exists \pi \in H,\ \pi(a) = b$;
  \item $\forall a \le b \in \mathbb{Q},\ \forall \pi \in H,\ \pi(b) = b \Rightarrow \pi(a) = a$;
  \item $\forall a < b < c \in \mathbb{Q},\ \exists \pi \in H,\ \pi(a) = a \ \text{and} \ \pi(c) < b$.
\end{enumerate}
\end{definition}

The existence of a KnightAutGroup is proved by a Baire-category / forcing argument:
one constructs a generic sequence of finite partial order-automorphisms of $\mathbb{Q}$ satisfying some carefully-chosen closed conditions and takes the group generated by the resulting generic automorphisms.

\begin{definition}[Generic KnightAutGroup]
  \label{def:knight-aut-group-forcing}
  \lean{knightAutGroup}
  \leanok
  \uses{def:knight-aut-group}
  The \emph{generic KnightAutGroup} $H$ is the subgroup of $\mathrm{Aut}(\mathbb{Q}, {<})$ generated by the generic automorphisms $\{\gamma_n \mid n \in \mathbb{N}\}$ produced by Baire category argument.
\end{definition}

Much of the meat of this formalization is behind Definition \ref{def:knight-aut-group-forcing}.

\begin{theorem}
  \label{thm:knight-aut-group}
  \uses{def:knight-aut-group-forcing}
  A KnightAutGroup exists.
\end{theorem}

\begin{proof}
Definition \ref{def:knight-aut-group-forcing}.
\end{proof}


\section{Blueprint}\label{sec:blueprint}

Given a KnightAutGroup $H$, the set $\{q \in \mathbb{Q} \mid q \le 1\}$ carries a natural blueprint structure: the multiplication $m \cdot n$ is defined as $h_n(m)$, where $h_n \in H$ is the unique element mapping $1 \mapsto n$.

\begin{definition}[Blueprint from a KnightAutGroup]
  \label{def:knight-aut-group-blueprint}
  \lean{KnightAutGroup.Blueprint}
  \leanok
  Given a KnightAutGroup $H$, its derived \emph{blueprint} $(\Omega_H, \le, \circ)$ is the subtype $\{q : \mathbb{Q} \mid q \le 1\}$ with multiplication $m \cdot n = g_n(m)$ and the usual order of the rationals.
\end{definition}

\begin{theorem}
  \label{thm:knight-aut-group-is-blueprint}
  \lean{KnightAutGroup.Blueprint.instBlueprint}
  \leanok
  \uses{def:blueprint, def:knight-aut-group-blueprint}
  Given a KnightAutGroup H, the derived blueprint $(\Omega_H, \le, \circ)$ does indeed satisfy the axioms of being a blueprint.
  In particular, a blueprint exists.
\end{theorem}

\begin{proof}
  \leanok
  This follows directly from the axioms.
\end{proof}


\section{Knight Models}\label{sec:knight-model}

\begin{definition}[Knight model from a blueprint]
  \label{def:knight-model-blueprint}
  \lean{Blueprint.Model}
  \leanok
  Given a blueprint $(\Omega, \circ, \le)$, the derived Knight model $\mathcal{K}_\Omega = (K, \preceq, f_n)_{n \in \Omega}$ is defined by $K = \{a \in \Omega \mid a < 1\}$ where $f_n(a) := n \circ a$ and $\preceq$ is just $\le$.
\end{definition}

\begin{theorem}
  \label{thm:blueprint-is-knight-model}
  \lean{Blueprint.Model.instKnightModel}
  \leanok
  Given a blueprint $(\Omega, \circ, \le)$ the derived Knight model $\mathcal{K}_\Omega$ is indeed a Knight model with blueprint $\Omega$.
\end{theorem}

\begin{proof}
  \leanok
  This follows directly from the axioms.
\end{proof}


\begin{definition}[Knight embeddings and isomorphisms]
  \label{def:knight-iso}
  \lean{KnightEmbedding, KnightIso, PartialKnightIso}
  \leanok
  \uses{def:knight-model}
  A \emph{Knight embedding} $\alpha \hookrightarrow \beta$ is an order embedding that commutes with all $f_n$. A \emph{Knight isomorphism} is a bijective Knight embedding. A \emph{partial Knight isomorphism} is a finitely-supported partial map that is order-preserving and commutes with all $f_n$ on its domain.
\end{definition}

\begin{theorem}[Countable uniqueness]
  \label{thm:countable-generic-iso}
  \lean{ModelIsoCondition.genericIso}
  \leanok
  \uses{def:knight-iso}
  Any two countable nonempty Knight models of the same blueprint $\Omega$ are isomorphic.
\end{theorem}

\begin{proof}
  \leanok
  \uses{def:knight-iso}
  Follows by a back-and-forth argument. The isomorphism is constructed via a Cohen-style generic filter meeting all necessary dense open sets of partial Knight isomorphisms.
\end{proof}

\begin{theorem}[Cardinality bound]
  \label{thm:knight-model-card-bound}
  \lean{KnightModel.knightModel_card_le_aleph_one}
  \leanok
  \uses{def:knight-model}
  Every Knight model has cardinality at most $\aleph_1$.
\end{theorem}

\begin{proof}
  \leanok
  \uses{def:knight-model}
  It is $\omega_1$-like linear order (every proper initial segment is countable).
\end{proof}


\section{Chains and Existence}\label{sec:chains}

The final step is to build, for a given blueprint, a Knight model of cardinality $\aleph_1$. This is done by taking a direct limit of a Zorn-maximal chain of Knight model embeddings indexed by countable ordinals.

\begin{definition}[Knight chain]
  \label{def:knight-chain}
  \lean{KnightChain}
  \leanok
  \uses{def:knight-model, def:knight-iso}
  A \emph{Knight chain} for blueprint $\Omega$ over a countable nonempty Knight model $K$ is a directed system of Knight embeddings $K \hookrightarrow K$, indexed by an initial segment of $\omega_1$, where every transition embedding factors through a fixed non-surjective self-embedding of $K$, ensuring the colimit strictly grows at each stage.
\end{definition}

\begin{theorem}[Uncountable colimit]
  \label{thm:chain-colimit-uncountable}
  \lean{KnightChain.not_countable_colimit}
  \leanok
  \uses{def:knight-chain}
  If the index domain of a Knight chain is uncountable, then its colimit is uncountable.
\end{theorem}

\begin{proof}
  \leanok
  \uses{def:knight-chain}
  The non-surjective transition embeddings ensure a genuinely new element is added at each successor stage; an uncountable index domain therefore produces uncountably many distinct elements in the colimit.
\end{proof}

\begin{theorem}[Full-domain chain]
  \label{thm:chain-full-domain}
  \lean{KnightChain.exists_fullDomain}
  \leanok
  \uses{def:knight-chain}
  By Zorn's lemma, there exists a Knight chain whose index domain is all of $\omega_1$.
\end{theorem}

\begin{proof}
  \leanok
  \uses{def:knight-chain}
  The set of Knight chains is partially ordered by extension. Every totally-ordered subfamily has an upper bound (union of the chain). A maximal element must have full domain $\omega_1$, since otherwise one could extend it by one more step.
\end{proof}

\begin{theorem}[Uncountable model existence]
  \label{thm:chain-uncountable-model}
  \lean{KnightChain.exists_uncountable_knightModel}
  \leanok
  \uses{def:knight-chain, thm:chain-full-domain, thm:chain-colimit-uncountable, thm:countable-generic-iso}
  For any blueprint $\Omega$, there exists an uncountable Knight model of type $\Omega$.
\end{theorem}

\begin{proof}
  \leanok
  \uses{def:knight-chain, thm:chain-full-domain, thm:chain-colimit-uncountable,
        thm:countable-generic-iso}
  Start with a countable Knight model.
  Apply Theorem~\ref{thm:chain-full-domain} to obtain a full-domain Knight chain; its colimit is a Knight model.
  By Theorem~\ref{thm:chain-colimit-uncountable} the colimit is uncountable (the domain $\omega_1$ is uncountable).
\end{proof}

\bibliographystyle{alpha}
\bibliography{refs}